\begingroup
\renewcommand{\arraystretch}{1.45}
\clearpage\fancyhead[L]{\small TRANSCENDENTAL FUNCTIONS}
\pdfbookmark[0]{PART II. TRANSCENDENTAL FUNCTIONS}{agm-pars2}
\SourceLeaf{381}
\begin{center}
\textbf{PART II.}\par
\textbf{ON THE TRANSCENDENTAL FUNCTIONS ARISING FROM THE DIFFERENTIATION}\par
\textbf{OF ARITHMETIC–GEOMETRIC MEANS.}
\end{center}
\par\begin{center}\rule{.15\linewidth}{.4pt}\end{center}
\Needspace{43\baselineskip}
\AGMArticle{9}
Let
\[
\left\{\begin{array}{llllll}
x,&x',&x'',&x'''\; .\; .\; .\; .&&\text{(I)}\\[6pt]
y,&y',&y'',&y'''\; .\; .\; .\; .&&\text{(II)}
\end{array}\right.
\]
be progressions formed in the same way as the progressions in Article 1, namely so that each term in
$\left\{\begin{array}{c}\mathrm{I}\\\mathrm{II}\end{array}\right\}$ is the
$\left\{\begin{array}{c}\text{arithmetic}\\\text{geometric}\end{array}\right\}$
mean between the two preceding terms, and thus $x^{\infty}=y^{\infty}=\mathrm{M}(x,y)$:
it is clear that all these quantities, and hence also their limit, are functions
of the two variables $x,y$, and change as soon as either or both of these
undergo a change. Here we shall speak only of infinitely small changes.
\Needspace{24\baselineskip}
It follows, therefore, that
\[
\begin{aligned}&\mathrm{d}x'=\tfrac12\mathrm{d}x+\tfrac12\mathrm{d}y,\\&
\mathrm{d}y'=\tfrac12\surd\frac{y}{x}. \mathrm{d}x
+\tfrac12\surd\frac{x}{y}. \mathrm{d}y\\&=\tfrac12\frac{y'}{x}\mathrm{d}x+\tfrac12\frac{y'}{y}\mathrm{d}y\end{aligned}
\]
and. similarly
\[
\begin{aligned}&\mathrm{d}x''=\tfrac12\mathrm{d}x'+\tfrac12\mathrm{d}y',\\&
\mathrm{d}y''=\tfrac12\frac{y''}{x'}. \mathrm{d}x'
+\tfrac12\frac{y''}{y'}\mathrm{d}y'\quad\text{etc.}\end{aligned}
\]
In this way the differentials of all the terms in I and II could be expressed through the differentials of
$x,y$ by substitutions; but the law of the progression would emerge in a very
obscure form. To bring it before our eyes as clearly as possible, we shall use the following
abbreviations: we shall write
\Needspace{20\baselineskip}
\SourceLeaf{382}
\[
\begin{array}{lll}
\dfrac{\mathrm{d}x}{x}+\dfrac{\mathrm{d}y}{y}=f,&
\dfrac{\mathrm{d}x'}{x'}+\dfrac{\mathrm{d}y'}{y'}=f',&
\dfrac{\mathrm{d}x''}{x''}+\dfrac{\mathrm{d}y''}{y''}=f''\quad\mathrm{etc.}\\[7pt]
\dfrac{\mathrm{d}x}{x}-\dfrac{\mathrm{d}y}{y}=g,&
\dfrac{\mathrm{d}x'}{x'}-\dfrac{\mathrm{d}y'}{y'}=g',&
\dfrac{\mathrm{d}x''}{x''}-\dfrac{\mathrm{d}y''}{y''}=g''\quad\mathrm{etc.}
\end{array}
\]
\Needspace{34\baselineskip}
whence it follows at once that
\[
\begin{aligned}
f'&=\tfrac12\left(\frac{\mathrm{d}x}{x'}+\frac{\mathrm{d}x}{x}+\frac{\mathrm{d}y}{x'}+\frac{\mathrm{d}y}{y}\right)\\&=f+\tfrac12\frac{\mathrm{d}x}{xx'}(x-x')+\tfrac12\frac{\mathrm{d}y}{yx'}(y-x')\\[7pt]
&=f+\tfrac12\frac{\mathrm{d}x}{xx'}(x-x')-\tfrac12\frac{\mathrm{d}y}{yx'}(x-x')\\&=f+\tfrac12 g\frac{x-x'}{x'}
\end{aligned}
\]
and in precisely the same way
\[
f''=f'+\tfrac12 g'\frac{x'-x''}{x''},\quad
f'''=f''+\tfrac12 g''\frac{x''-x'''}{x'''}\quad\mathrm{etc.}
\]
\Needspace{27\baselineskip}
Similarly,
\[
\begin{aligned}&g'=\tfrac12\left(\frac{\mathrm{d}x}{x'}-\frac{\mathrm{d}x}{x}+\frac{\mathrm{d}y}{x'}-\frac{\mathrm{d}y}{y}\right)\\&=\tfrac12\frac{\mathrm{d}x}{xx'}(x-x')+\tfrac12\frac{\mathrm{d}y}{yx'}(y-x')\\&=\tfrac12 g\left(\frac{x-x'}{x'}\right)\end{aligned}
\]
and in the same way \(g''=\tfrac12 g'\frac{x'-x''}{x''},\quad
g'''=\tfrac12 g''\frac{x''-x'''}{x'''}\) etc.
It is also clear from this that
\(f'=f+g',\ f''=f+g'+g'',\ f'''=f+g'+g''+g'''\) etc.
\Needspace{25\baselineskip}
It is readily seen that the sequence
\(\frac{g'}{g},\frac{g''}{g},\frac{g'''}{g}\) etc.
decreases extremely rapidly below every bound; hence we shall have, by an infinite series converging extremely rapidly,
\[
\begin{aligned}&f^\infty=f+g\Bigl\{\tfrac12\frac{x-x'}{x'}
+\tfrac14\frac{x-x'}{x'}\cdot\frac{x'-x''}{x''}\\&\qquad+\tfrac18\frac{x-x'}{x'}\cdot\frac{x'-x''}{x''}\cdot\frac{x''-x'''}{x'''}
+\mathrm{etc.}\Bigr\},\\&g^\infty=0\end{aligned}
\]\Needspace{24\baselineskip}

hence
\[
\begin{aligned}
\frac{2\mathrm{d}x^\infty}{x^\infty}
&=\frac{2\mathrm{d}y^\infty}{y^\infty}
=\frac{2\mathrm{d}\mathrm{M}(x,y)}{\mathrm{M}(x,y)}\\[7pt]
&=\frac{\mathrm{d}x}{x}+\frac{\mathrm{d}y}{y}\\&\quad+\left\{\frac{\mathrm{d}x}{x}-\frac{\mathrm{d}y}{y}\right\}\times
\left\{\tfrac12\frac{x-x'}{x}+\tfrac14\frac{x-x'}{x'}\cdot\frac{x'-x''}{x''}+\mathrm{etc.}\right\}
\end{aligned}
\]
\Needspace{28\baselineskip}
or
\[
\begin{aligned}
&\mathrm{d}\mathrm{M}(x,y)=\mathrm{M}(x,y)\times\Bigl\{\\&\quad\frac{\mathrm{d}x}{x}\times\left(\tfrac12+\tfrac14\frac{x-x'}{x'}+\tfrac18\frac{x-x'}{x'}\cdot\frac{x'-x''}{x''}+\mathrm{etc.}\right)\\&\quad+\frac{\mathrm{d}y}{y}\times\left(\tfrac12-\tfrac14\frac{x-x'}{x'}-\tfrac18\frac{x-x'}{x'}\cdot\frac{x'-x''}{x''}-\mathrm{etc.}\right)\Bigr\}
\end{aligned}
\]
\Needspace{39\baselineskip}
\par We obtain a somewhat more convenient calculation as follows: it follows that
\[
\begin{aligned}&\frac{x-x'}{x'}=\frac{x-y}{x+y}=\frac{(x-y)^2}{xx-yy}\\&=\frac{(x+y)^2-4xy}{xx-yy}=4\frac{x'x'-y'y'}{xx-yy}\text{\textsuperscript{*)}}\end{aligned}
\]\footnotetext{\csname HistoricalNoteText30\endcsname}
whence follows
\[
\begin{aligned}
&\mathrm{d}\mathrm{M}(x,y)=\frac{\mathrm{M}(x,y)}{2(xx-yy)}\times\Bigl\{\\&\quad\frac{\mathrm{d}x}{x}\times\bigl(xx-yy+2(x'x'-y'y')\\&\qquad+4(x''x''-y''y'')+8(x'''x'''-y'''y''').\;.\;.\bigr)\\&\quad+\frac{\mathrm{d}y}{y}\times\bigl(xx-yy-2(x'x'-y'y')\\&\qquad-4(x''x''-y''y'')-8(x'''x'''-y'''y''').\;.\;.\bigr)\Bigr\}
\end{aligned}
\]
\Needspace{26\baselineskip}
\SourceLeaf{383}
\par 
To make the calculation as easy as possible by means of logarithms, we append the following formulas, which follow directly from the algorithm of progressions I. II:
\[
(x'x'-y'y')=\frac{(xx-yy)^2}{16x'x'},\quad
x''x''-y''y''=\frac{(x'x'-y'y')^2}{16x''x''}\quad\mathrm{etc.}
\]
these also show at once how rapidly our series must converge.

\par If preferred, the following formulas may also be used:
\[
x'x'-y'y'=\tfrac14(x-y)^2,\quad
x''x''-y''y''=\tfrac14(x'-y')^2\quad\mathrm{etc.}
\]

\clearpage
\endgroup
